                          %——— PARTE N ———%                           
%%%≠≠≠≠≠≠≠≠≠≠≠≠≠≠≠≠≠≠≠≠≠≠≠≠≠≠≠≠≠≠≠≠≠≠≠≠≠≠≠≠≠≠≠≠≠≠≠≠≠≠≠≠≠≠≠≠≠≠≠≠≠≠≠≠≠≠≠≠≠≠%
% Modello per la struttura "per parte": copiare in N_Titolo.tex, una copia per parte,
% e richiamarle da main.tex tra introduzione e conclusioni. Questo file non viene compilato.
\section{TITOLO DELLA PARTE}
% Una frase: cosa si vuole misurare o verificare.

%%%——————————————————————————————————————————————————————————————————————%
\paragraph{Scelte di progetto}
% Solo se la consegna lascia liberi elementi o valori: cosa si è scelto e perché
% (vincoli degli strumenti, limiti dei componenti, richieste della consegna).
% Altrimenti togliere il paragrafo.
%——————————————————————————————————————————————————————————————————————%%%

%%%——————————————————————————————————————————————————————————————————————%
\paragraph{Circuito e procedura}
% Schema o foto (images/), valori misurati dei componenti, cosa si è letto
% e con quale strumento, come si è stimata l'incertezza.
%——————————————————————————————————————————————————————————————————————%%%

%%%——————————————————————————————————————————————————————————————————————%
\paragraph{Misure}
% Tabella con i valori di data/ (vedi snippets.tex).
%——————————————————————————————————————————————————————————————————————%%%

%%%——————————————————————————————————————————————————————————————————————%
\paragraph{Risultati}
% Grafico da plots/, risultati con \res{...}, confronto con l'atteso, commento breve.
%——————————————————————————————————————————————————————————————————————%%%
